
Machine learning benchmarks serve as primary instruments for measuring research progress across computer vision (CV), natural language processing (NLP), and tabular domains. As benchmarks accumulate submissions over time, score distributions may narrow, reducing the ability of the benchmark to discriminate between genuinely different models. This phenomenon — referred to here as \textit{score compression} — has been noted qualitatively for benchmarks such as MMLU and CIFAR-10, but its prevalence, temporal dynamics, and causal structure have not been measured systematically across domains.

The broader problem of benchmark saturation is recognized in the ML research community. Prior work has characterized overfitting to benchmark test sets in CV using retesting studies (Roelofs et al., 2019; Recht et al., 2019), provided saturation indices for LLM benchmarks (Polo et al., 2026), and developed contamination detection tools for NLP (Chan et al., 2024). However, no cross-domain analysis has established whether submission accumulation \textit{temporally precedes} compression — a prerequisite for any prospective monitoring framework — nor has compression prevalence been measured at field scale across CV, NLP, and tabular benchmarks simultaneously.

This paper describes the Benchmark-Calibrated Health Score (BCBHS) framework, which was designed to address this gap. The framework's original goal was to build a Cox proportional hazards model that predicts time-to-discriminative-collapse for individual benchmarks, using domain-specific health estimators H_d as time-varying covariates. In this work, two precondition experiments were executed (H-E1 and H-M1) and one intermediate experiment was attempted (H-M2); the survival model itself (H-M3 and H-M4) was not executed due to a collapse event operationalization failure identified during H-M2.

The main empirical contributions of this work are as follows:

\begin{enumerate}
\item A measurement of benchmark compression prevalence: 31.1\% of 466 qualifying benchmarks exhibit at least one compression event over a 7-year leaderboard panel (2018–2025).
\item Granger causality evidence that submission count accumulation temporally precedes score variance compression at lag 2 (p = 1.854 × 10⁻⁵), with reverse causality not confirmed.
\item Discriminability results for domain-specific H_d signals on synthetic panels calibrated to real data (|d| > 5 in all domains), along with NLP cross-sectional AUC = 0.857 on real benchmark data.
\item Identification and root-cause analysis of a structural incompatibility in the pre-registered collapse operationalization, with a proposed resolution path.
\end{enumerate}

This paper does not claim to provide a validated prospective prediction system. The C-index and lead-time components of the BCBHS hypothesis were not reached due to the pipeline failure at H-M2.

